\beginsong{Când ajungi să-ți fii tu însuți o povară}
\beginverse
Când ajungi să-ți fii tu însuți o povară,
Când dureri de ne-nțeles te-apasă greu,
Când pe crengi sunt numai flori de primăvară,
Iar în tine frunze moarte cad mereu,
Când spre tainice urzeli ce se destramă
Rătăcești ca un pescar de amăgiri,
N-auzi tu un glas de frate ce te cheamă?
N-auzi tu suspinul veșnicei iubiri?
\endverse

\beginchorus
Întoarce-te la Domnul, așa cum ești,
Cu zdrențe-n suflet și poveri!
Întoarce-te la Domnul, un cuget nou
Și-o altă inimă să-I ceri!
Întoarce-te la Domnul, la pieptul Lui,
Să înțelegi cât te-a iubit.
Pe rănile eterne, vei plânge fericit,
Întoarce-te la Domnul, copil regăsit.
\endchorus

\beginverse
Ca-ntr-un joc, legându-ți ochii cu năframa,
Tu alergi să prinzi minciuna unui vis,
Dar oprește-te, prietene, ia seama!
Pașii tăi se-ndreaptă-n goană spre abis!
În zadar pândești năluca fericirii,
În zadar apuci cu pumnii amândoi,
Vei plăti din nou tribut dezamăgirii,
Obosit, cu ochii plânși, cu pumnii goi!
\endverse
\endsong